\documentclass[dvipdfmx]{jsarticle}
\usepackage[dvipdfmx]{xcolor}
\usepackage{amsmath}
\usepackage{amssymb}
\usepackage{amsthm}
\usepackage{mathtools}
\usepackage{bm}
\usepackage[margin=15truemm]{geometry}
\usepackage{titlesec}
\usepackage[dvipdfmx]{graphicx}
\usepackage{tikz}
\usepackage{ascmac} 
\usepackage{enumitem}
\usepackage{thmbox}
\usepackage{pgfplots}
\usepackage{physics}
\usepackage{float}
\usepackage{quantikz}
\usepackage{adjustbox}

\title{\vspace{-2cm}量子フーリエ変換}
\author{河田直也}
\date{\today}
\raggedbottom
\titlespacing*{\section}{0pt}{*0}{*0}
\begin{document}
\maketitle
\section{量子フーリエ変換とは}
離散フーリエ変換は以下のように定義される。通常の文脈では指数部分の符号はマイナスであるが、量子フーリエ変換との接続を考えて符号をプラスにしている。通常の逆フーリエ変換ととらえることができる。
\begin{equation}
    y_k := \frac{1}{\sqrt{N}} \sum_{j=0}^{N-1} x_j e^{2\pi ijk/N}\label{eq:dft}
\end{equation}
$x_0,\cdots, x_{N-1}$の入力に対して、$y_0,\cdots, y_{N-1}$の出力をするような変換になっていることがわかる。量子フーリエ変換は離散フーリエ変換を量子の表現で置き換えたものである。\\
量子フーリエ変換は$2^n$次元のHilbert空間$\mathcal{H}$上の$N$個の正規直交基底$\ket{0},\cdots, \ket{N-1}$を別の$N$個の正規直交基底$\left\{\ket{\tilde{k}}\right\}_{k=0}^{N-1}$に写すユニタリ変換である。
\begin{equation}
    \ket{j} \rightarrow \frac{1}{\sqrt{N}} \sum_{k=0}^{N-1} e^{2\pi ijk/N} \ket{k} \label{eq:qft_basis}
\end{equation}
同様にして$N$個の基底の線形結合であらわされる任意の状態に対する変換として
\begin{align}
    \sum_{j=0}^{N-1} x_j \ket{j} &\rightarrow \sum_{j=0}^{N-1}x_j\left(\frac{1}{\sqrt{N}}\sum_{k=0}^{N-1} e^{2\pi ijk/N} \ket{k}\right) \notag \\
    &= \frac{1}{\sqrt{N}} \sum_{k=0}^{N-1} \left(\sum_{j=0}^{N-1} x_j e^{2\pi ijk/N}\right) \ket{k}\notag \\
    &= \sum_{k=0}^{N-1} y_k \ket{k}\quad(\because \text{式\ref{eq:dft}})\label{eq:qft}
\end{align}
と書けることもわかる。ここで$y_k$は振幅$x_j$の離散フーリエ変換での出力である。
\vspace{5mm}
\section{ユニタリ性}
量子フーリエ変換が量子アルゴリズムであるためにはユニタリである必要がある。量子フーリエ変換がユニタリであることは以下の2通りの手法で示すことができる。
\begin{enumerate}[label = (\roman*)]
    \item 量子フーリエ変換$\hat{F}$の行列表示$F$がユニタリ行列であることを示す
    \item 量子フーリエ変換を行う量子回路をユニタリな量子ゲートのみを用いて構成できることを示す
\end{enumerate}
量子回路の構成は必然的に述べるため、ここでは(i)の方法での証明を与える。
\begin{proof}
    量子フーリエ変換$\hat{F}$を正規直交基底$\ket{0},\cdots,\ket{N-1}$において行列表示する。行列$F$の成分$F_jk=\bra{j}\hat{F}\ket{k}$は以下のとおりである。
    \begin{align}
        F_{jk} &= \bra{j}\hat{F}\ket{k}\notag\\
        &=\bra{j}\left(\frac{1}{\sqrt{N}}\sum_{l=0}^{N-1}e^{2\pi ikl/N}\ket{l}\right)\notag\\
        &= \frac{1}{\sqrt{N}}e^{2\pi ijk/N}\quad(\because \text{正規直交基底})
    \end{align}
    次に、$F^\dagger F$の成分を計算して$F^\dagger F = I$が成り立つことを示す。
    \begin{align}
        (F^\dagger F)_{jk} &= \sum_{l=0}^{N-1} (F^\dagger)_{jl} F_{lk} \notag \\
        &= \sum_{l=0}^{N-1} \overline{F_{lj}} F_{lk} \quad (\because \text{$\dagger$は共役転置を意味する})\notag \\
        &= \sum_{l=0}^{N-1} \left(\frac{1}{\sqrt{N}}e^{-2\pi ilj/N}\right) \left(\frac{1}{\sqrt{N}}e^{2\pi ilk/N}\right) \notag \\
        &= \frac{1}{N} \sum_{l=0}^{N-1} e^{2\pi il(k-j)/N} \notag \\
        &= \delta_{jk} \quad(\because \text{等比数列の和の公式を用いる})
    \end{align}
    よって$\hat{F}$がユニタリであることが示された。
\end{proof}
これはフーリエ変換をした後に逆フーリエ変換を行うと元に戻るという性質の表れといえる。量子フーリエ変換がユニタリであることが示されたので、ユニタリな量子ゲートを用いてアルゴリズムの実装が可能ということが分かった。
\vspace{5mm}
\section{積表示}
$N=2^n$とする。$n$量子ビット量子コンピュータでの計算基底が$\ket{0},\cdots, \ket{2^n-1}$の計$2^n$個であるためである。ここで、
\[
    j = j_1 2^{n-1} + j_2 2^{n-2} + \cdots +j_n2^0
\]
と表される$j$は2進数表示すると
\[
    j = j_1j_2\cdots j_n {}_{(2)}
\]
となる。さらに、$1$より小さい数についても同様に考えることができ、
\[
    j = \frac{j_l}{2}+\frac{j_{l+1}}{4}+\cdots + \frac{j_m}{2^{m-l+1}}
\]
と表される$j$は2進数表示すれば
\[
    j = 0.j_l j_{l+1}\cdots j_m {}_{(2)}
\]
となる。この2進数表示を用いると量子フーリエ変換の出力は以下のように表すことができる。（\textbf{積表示}）
\begin{equation}
    \ket{j_1\cdots j_n}\rightarrow \frac{\left(\ket{0}+e^{2\pi i0.j_n {}_{(2)}}\ket{1}\right)\left(\ket{0}+e^{2\pi i0.j_{n-1}j_n {}_{(2)}}\ket{1}\right)\cdots \left(\ket{0}+e^{2\pi i0.j_1\cdots j_n {}_{(2)}}\ket{1}\right)}{2^{n/2}}
\end{equation}
この積表示は量子フーリエ変換の定義（式\ref{eq:qft_basis}）から簡単な計算によって求めることができる。以下にその計算を示す。
\begin{align}
    \ket{j} &\rightarrow \frac{1}{2^{n/2}} \sum_{k=0}^{2^n-1} e^{2\pi ijk/2^n} \ket{k}\quad(\because \text{式\ref{eq:qft_basis}において$N=2^n$の場合})\notag\\
     & = \frac{1}{2^{n/2}} \sum_{k_1=0}^1 \cdots \sum_{k_n=0}^1 e^{2\pi ij(k_1 2^{n-1} + \cdots + k_n)/2^n} \ket{k_1\cdots k_n}\quad(\because \text{$k$を2進数表示})\notag\\
     &= \frac{1}{2^{n/2}} \sum_{k_1=0}^1 \cdots \sum_{k_n=0}^1 e^{2\pi ij\left(\sum_{l=1}^n k_l 2^{-l}\right)} \ket{k_1\cdots k_n}\quad(\text{$1$より小さい数の2進数表示を思い出す})\notag \\
     &=\frac{1}{2^{n/2}} \sum_{k_1=0}^1 \cdots \sum_{k_n=0}^1 \bigotimes_{l=1}^n e^{2\pi ijk_l 2^{-l}}\ket{k_l} \notag \\
     &= \frac{1}{2^{n/2}} \bigotimes_{l=1}^n\left(\sum_{k_l=0}^1 e^{2\pi i j k_l 2^{-l}}\ket{k_l}\right)\quad(\text{テンソル積と和の交換})\notag\\
     &= \frac{1}{2^{n/2}} \bigotimes_{l=1}^n \left(\ket{0} + e^{2\pi ij2^{-l}}\ket{1}\right)\notag \\
     &=\frac{\left(\ket{0}+e^{2\pi i0.j_n {}_{(2)}}\ket{1}\right)\left(\ket{0}+e^{2\pi i0.j_{n-1}j_n {}_{(2)}}\ket{1}\right)\cdots \left(\ket{0}+e^{2\pi i0.j_1\cdots j_n {}_{(2)}}\ket{1}\right)}{2^{n/2}}
\end{align}
\vspace{5mm}
\section{量子回路}
次に、量子フーリエ変換を実装するための量子回路を構成する。この構成に際して、積表示が重要な役割を果たす。
\begin{figure}[H]
    \centering
\begin{adjustbox}{width=1.0\textwidth}
\begin{quantikz}
    \lstick{$\ket{j_1}$} & \gate{H} & \gate{R_2} \qw & \dots & \gate{R_{n-1}} \qw & \gate{R_n} \qw & \qw & \qw & \qw &\qw&\qw&\qw&\qw &\qw & \rstick{$\frac{1}{\sqrt{2}}(\ket{0} + e^{2\pi i 0.j_1 \dots j_n} \ket{1})$}\\
    \lstick{$\ket{j_2}$} & \qw & \ctrl{-1} & \dots & \qw & \qw & \gate{H} & \dots & \gate{R_{n-2}} & \gate{R_{n-1}} &\dots & \qw &\qw &\qw &   \rstick{$\frac{1}{\sqrt{2}}(\ket{0} + e^{2\pi i 0.j_2 \dots j_n} \ket{1})$} \\
    \vdots &\setwiretype{n} & & \vdots & & & & & & & \vdots & & & & &\vdots \\
    \lstick{$\ket{j_{n-1}}$} & \qw & \qw & \qw & \ctrl{-3} & \qw & \qw & \qw & \ctrl{-2} & \qw & \cdots &  \gate{H} &\gate{R_2}&\qw & \rstick{$\frac{1}{\sqrt{2}}(\ket{0} + e^{2\pi i 0.j_{n-1} j_n} \ket{1})$} \\
    \lstick{$\ket{j_n}$} & \qw & \qw & \qw & \qw & \ctrl{-4} & \qw & \qw &\qw &  \ctrl{-3} &\cdots &\qw & \ctrl{-1}&\gate{H}&\rstick{$\frac{1}{\sqrt{2}}(\ket{0} + e^{2\pi i 0.j_n} \ket{1})$}\label{eq:product}
\end{quantikz}
\end{adjustbox}
    \caption{量子フーリエ変換の量子回路（SWAPゲートを省略）}
\end{figure}
以上が$n$量子ビットにおける量子フーリエ変換の量子回路である。回路に登場する$R_k$ゲートの行列表示は以下のとおりである。
\begin{equation}
    R_k := 
    \begin{pmatrix}
        1 & 0 \\
        0 & e^{2\pi i/2^k}\\
    \end{pmatrix}
\end{equation}
つまり、定義より$R_k$は位相を$\pi/2^k$シフトさせる量子ゲートであることがわかる。また、$k=2$のとき$S$ゲート、$k=3$のとき$T$ゲートであることもわかる。\\
それでは、具体的に入力$\ket{j_1j_2\cdots j_n}$が量子回路を通る過程でどう変化するかを追うことにする。
\begin{enumerate}[label = (\arabic*)]
    \item \textbf{量子ビット1への操作}\\
    状態の入力に対して、まず量子ビット1に$H,R_2,\cdots, R_{n-1},R_n$の順にゲート操作を行うことになる。$H$ゲートは
    \[
        \ket{0}\rightarrow \frac{\ket{0}+\ket{1}}{\sqrt{2}},\quad\ket{1}\rightarrow \frac{\ket{0}-\ket{1}}{\sqrt{2}}
    \]
    と状態を変化させるゲートであり、これは入力$\ket{j_1}(j_1=0\:\text{or}\:1)$に対して
    \[
        \ket{j_1}\rightarrow \frac{1}{2^{1/2}}\left(\ket{0}+e^{2\pi i0.j_1}\ket{1}\right)\label{eq:H}
    \]
    と統一的に表現できる。$j_1=0$なら$e^{2\pi i 0.0}= 1$、$j_1=1$なら$e^{2\pi i0.1}=-1$なので$H$ゲートによる変換になっている。\\
    次に、制御$R_2$ゲートをかける。制御$R_2$ゲートは量子ビット2が0なら恒等変換、1なら$R_2$ゲートをかけるゲートである。$R_2$ゲートによる変換を具体的に記すと
    \[
        \ket{0}\rightarrow \ket{0},\quad \ket{1}\rightarrow \frac{\ket{0}+i\ket{1}}{\sqrt{2}}
    \]
    となり、式\ref{eq:H}の状態に対して制御$R_2$ゲートをかけると
    \[
        \frac{1}{2^{1/2}}\left(\ket{0}+e^{2\pi i0.j_1j_2}\ket{1}\right)
    \]
    さらに制御$R_3$、制御$R_4$、と繰り返して制御$R_n$ゲートまでかけると、状態は結果的に$\ket{j_1}$から以下のように変わる。
    \[
        \ket{j_1}\rightarrow \frac{1}{2^{1/2}}\left(\ket{0}+e^{2\pi i0.j_1j_2\cdots j_n}\ket{1}\right)
    \]
    \item \textbf{量子ビット2への操作}\\
    量子ビット1への操作が終わると量子ビット2への操作を行うことになる。量子ビット2へは$H,R_2,\cdots,R_{n-1}$の順にゲート操作を行う。量子ビット1への操作と同様に考えて、入力$\ket{j_2}(j_2=0\:\text{or}\:1$)に対して
    \[  
        \ket{j_2} \rightarrow \frac{1}{2^{1/2}}\left(\ket{0}+e^{2\pi i0.j_2\cdots j_n}\ket{1}\right)
    \]
    となる。
    \item \textbf{量子ビットnまで操作を繰り返す}\\
    量子ビット1,量子ビット2にした操作を量子ビットnまで繰り返すと以下の状態が得られることがわかる。
    \begin{equation}
        \frac{1}{2^{n/2}}\left(\ket{0}+e^{2\pi i0.j_1j_2\cdots j_n}\ket{1}\right)\left(\ket{0}+e^{2\pi i0.j_2\cdots j_n}\ket{1}\right)\cdots \left(\ket{0}+e^{2\pi i0.j_n}\ket{1}\right)\label{eq:beforeswap}
    \end{equation}
    \item \textbf{SWAPゲートをかける}\\
    量子フーリエ変換における積表示(式\ref{eq:product})と量子回路を通った出力（式\ref{eq:beforeswap}）を比較するとビットの並び順が逆になっていることがわかる。そのため、出力された状態にSWAPゲートをかけることで得たい出力を得る。
\end{enumerate}
以上の流れで量子フーリエ変換の量子回路になっていることが確かめられた。量子フーリエ変換を実行できる量子回路を、ユニタリゲートのみで構成することができたことも量子フーリエ変換がユニタリ変換であることを示している。
\vspace{5mm}
\section{計算量}
最後に量子フーリエ変換における計算量と古典コンピュータでの計算量を比較する。
\begin{description}
    \item[量子]量子コンピュータの計算量を必要な量子ゲート数で見積もる。\\
    $N=2^n$において必要な量子ゲート数は量子ビットnまで操作を繰り返す段階にて、
    \[
        n+(n-1)+\cdots+1 = \frac{n(n+1)}{2}\text{個}
    \]
    さらに、SWAPゲートは$[n]/2$個必要で、SWAPゲート1個がCNOTゲート3個で実装されることを考えて
    \[  
        \frac{3n}{2}\text{個}
    \]
    である。以上より、量子フーリエ変換における計算量は$\mathcal{O}(n^2)$と見積もることができる。
    \item[古典]
    離散フーリエ変換の古典コンピュータを用いたアルゴリズムに高速フーリエ変換（FFT）がある。高速フーリエ変換の計算量は$\mathcal{O}(N\log N)$であり、$N=2^n$を代入すれば$\mathcal{O}(n 2^n)$といえる。
\end{description}
量子フーリエ変換と高速フーリエ変換の計算量を比較すると指数関数的に量子フーリエ変換の方が少ない計算量で離散フーリエ変換が行えることがわかる。量子フーリエ変換自体は驚くべき結果をもたらすようなものではないが、応用先に量子位相推定アルゴリズムやそれを利用したShorのアルゴリズム、離散対数問題がある。
\end{document}